% Recovered verbatim from chapters_31_35_with_part/chapters_31_35_with_part.tex
% Canonical chapter 34
\chapter{White-State Cycling and the Expansion History}

\section{Imprints from Previous Cycles}

In a cyclic MK--43 Ultra cosmology, each bounce passes through a white-state core configuration $\Wcore$.
Residual structures in the ER network and chromatic fields can survive in encoded form, influencing the initial conditions of the next cycle.
This may leave subtle imprints in the expansion history and structure formation patterns.

\section{Epoch Transitions Driven by Phase Thermodynamics}

The thermodynamic behaviour of the phase field and ER foam---as developed in Part~V---drives transitions between cosmological epochs:
\begin{itemize}
  \item stiff or kinetic-dominated phases,
  \item radiation-like and matter-like eras,
  \item dark-energy-like accelerated expansions.
\end{itemize}
These transitions may admit a unified description in terms of phase-stability bands and their crossings.

\section{Bounce Signatures in Observables}

Signatures of a preceding cycle's bounce may include:
\begin{itemize}
  \item anomalies at the largest observable scales in the CMB,
  \item features in the matter power spectrum,
  \item relic correlations in ER connectivity at super-horizon scales.
\end{itemize}
While such effects are subtle, they offer a route to test cyclic aspects of MK--43 Ultra.

